\subsection{Subsets and Power Sets}

\begin{definition} \label{def:proper-subsets-singletons-and-nonempty} \hlt{Basic Subset Terminology}\\
Let $A,B$ be sets.
\begin{enumerate}[label=(\roman*)]
\setlength{\itemsep}{0pt}
\item \hlt{Subset}: $A\subseteq B$ if and only if $x\in A$ implies $x\in B$ for every object $x$.
\item \hlt{Proper subset}: $A\subsetneq B$ iff $A\subseteq B$ and $A\neq B$.
\item \hlt{Singleton}: $A=\{x\}$ for some object $x$.
\item \hlt{Nonempty}: $A\neq\emptyset$.
\end{enumerate}
\end{definition}

\begin{assumption} \label{assump:set-extensionality} \hlt{Set Extensionality}\\
Throughout these notes, set equality is governed by extensionality.\\
For sets $A$ and $B$, $A=B$ if and only if $A\subseteq B$ and $B\subseteq A$.
\end{assumption}

\begin{proofmethod} \label{meth:proof-by-containment} \hlt{Proof by Two Containments}\\
To prove equality of sets $A$ and $B$, discharge the following obligations.
\begin{enumerate}[label=(\roman*)]
\setlength{\itemsep}{0pt}
\item Fix an arbitrary $x\in A$ and prove $x\in B$; this establishes $A\subseteq B$.
\item Fix an arbitrary $x\in B$ and prove $x\in A$; this establishes $B\subseteq A$.
\item Invoke Assumption \ref{assump:set-extensionality} to conclude $A=B$.
\end{enumerate}
If both sets have explicit membership criteria, the two inclusions may be proved together.\\
Fix an arbitrary object $x$ and prove $x\in A\Longleftrightarrow x\in B$.
\end{proofmethod}

\begin{definition} \label{def:power-set} \hlt{Power Set}\\
For a set $X$, its \hlt{power set} is $\mathcal{P}(X)\defeq\{S:S\subseteq X\}$.
The two smallest cases are as follows.
\begin{enumerate}[label=(\roman*)]
\setlength{\itemsep}{0pt}
\item The empty set has only itself as a subset, so $\mathcal{P}(\emptyset)=\{\emptyset\}$.
\item If $X=\{a,b\}$ with $a\neq b$, then $\mathcal{P}(X)=\{\emptyset,\{a\},\{b\},X\}$.
\end{enumerate}
\end{definition}
